\documentclass[aspectratio=169,11pt]{beamer}

\usetheme{default}
\usefonttheme{professionalfonts}
\setbeamertemplate{navigation symbols}{}
\setbeamertemplate{footline}[frame number]
\setbeamertemplate{itemize items}[circle]

\usepackage{amsmath,amssymb}
\usepackage{graphicx}
\usepackage{booktabs}
\usepackage{xcolor}
\usepackage{ragged2e}

\definecolor{HeatBlue}{RGB}{20,70,150}
\definecolor{HeatRed}{RGB}{205,45,35}
\setbeamercolor{structure}{fg=HeatBlue}
\setbeamercolor{frametitle}{fg=HeatBlue}

\title{HeatMapSignals}
\subtitle{Heat Diffusion as a 2D Signal Processing Problem}
\author{Farhan Anjum \and Naeem Shuvo}
\institute{Signals and Linear Systems}
\date{September 2026}

\newcommand{\img}[2][0.8\textwidth]{\includegraphics[width=#1,height=0.72\textheight,keepaspectratio]{#2}}

\begin{document}

% ------------------------------------------------
\begin{frame}
  \titlepage
\end{frame}

% ------------------------------------------------
\begin{frame}{Heat as a 2D Spatial Signal}
\centering
\img[0.82\textwidth]{heat_diffusion_over_time.png}

\vspace{1mm}
\begin{columns}[T,onlytextwidth]
\column{0.32\textwidth}
\small
\begin{itemize}
  \item The temperature field is a sampled 2D signal:
  \[
    T(x,y,t)
  \]
\end{itemize}

\column{0.32\textwidth}
\small
\begin{itemize}
  \item Each grid cell stores the temperature at that spatial location.
\end{itemize}

\column{0.36\textwidth}
\small
\begin{itemize}
  \item In each timestep we model:
  \begin{enumerate}
    \item Add heat
    \item Diffuse heat
    \item Cool toward ambient
    \item Boundary heat loss
  \end{enumerate}
\end{itemize}
\end{columns}
\end{frame}

% ------------------------------------------------
\begin{frame}{Thermodynamic Decay \& Boundary Loss}
\begin{columns}[T,onlytextwidth]
\column{0.5\textwidth}
\textbf{Temporal Cooling (Decay)}
\vspace{2mm}
\small
\begin{itemize}
  \item Heat dissipates into the ambient environment over time.
  \item Follows Newton's Law of Cooling (exponential decay):
  \[
    \boxed{T_{new} = T_{amb} + (T - T_{amb}) e^{-\lambda \Delta t}}
  \]
  \item $\lambda$: Decay rate
\end{itemize}

\column{0.5\textwidth}
\textbf{Boundary Loss}
\vspace{2mm}
\small
\begin{itemize}
  \item Heat escapes from the edges of the physical simulation grid.
  \item Applied to boundary pixels (Robin/absorbing boundary condition):
  \[
    \boxed{T_{edge} \leftarrow T_{edge} \times (1 - \gamma)}
  \]
  \item $\gamma$: Boundary loss rate
\end{itemize}
\end{columns}
\end{frame}

% ------------------------------------------------
\begin{frame}{Why a Gaussian Kernel?}
\centering
\begin{columns}[c,onlytextwidth]
\column{0.49\textwidth}
\centering
\img[0.98\textwidth]{heat_kernel_graph.png}

\column{0.49\textwidth}
\centering
\img[0.98\textwidth]{heat_kernel_matrix.png}
\end{columns}

\vspace{-2mm}
\begin{center}
\small
The fundamental solution of the 2D heat equation is Gaussian:
\[
\boxed{h(x,y,t)=\frac{1}{4\pi Dt}\exp\left(-\frac{x^2+y^2}{4Dt}\right)}
\]
\end{center}
\end{frame}

% ------------------------------------------------
\begin{frame}{Convolution: How Heat Spreads}
\begin{columns}[c,onlytextwidth]
\column{0.67\textwidth}
\centering
\img[1.0\textwidth]{kernel_sliding_convolution.png}

\column{0.33\textwidth}
\Large
\begin{itemize}
  \item The Gaussian kernel describes how heat at one location contributes to nearby locations.
  \item Sliding the kernel over the entire spatial field performs 2D spatial convolution.
\end{itemize}

\[
\boxed{S*h}
\]
\end{columns}
\end{frame}

% ------------------------------------------------
\begin{frame}{Why Fourier Transform?}
\centering
\img[0.92\textwidth]{spatial_frequency_domain.png}

\vspace{-2mm}
\begin{center}
\large
\textbf{Spatial domain} $\;\xrightarrow{\text{2D FFT}}\;$ \textbf{Frequency domain}
\end{center}
\end{frame}

% ------------------------------------------------
\begin{frame}{Direct Convolution vs. FFT Convolution}
\centering
\img[0.92\textwidth]{direct_vs_fft.png}

\vspace{1mm}
\begin{columns}[T,onlytextwidth]
\column{0.48\textwidth}
\small
\textbf{For our 100$\times$100 grid:}

\vspace{1mm}
\begin{tabular}{ll}
Grid cells & $N=10,000$\\
Kernel & $10\times10$\\
Kernel elements & $K=100$\\
\end{tabular}

\vspace{2mm}
\textbf{Direct spatial convolution}
\[
10,000\times100=\boxed{1,000,000}
\]
multiplications per frame.

\column{0.52\textwidth}
\small
\textbf{FFT-based convolution}
\[
N\log_2N
=10,000(13.28)
\approx\boxed{132,800}
\]
operations per frame.

\vspace{1mm}
The source PDF reports the FFT approach as \textbf{$>50\times$ faster} in the live app.\textsuperscript{[1]}
\end{columns}
\end{frame}

% ------------------------------------------------
\begin{frame}{The Convolution Theorem}
\centering
\vspace{5mm}
\Large
\[
\boxed{\mathcal{F}\{f*g\}=F(u,v)\,G(u,v)}
\]

\vspace{10mm}
\large
\textbf{Temperature / Source}
\quad $\rightarrow$ \quad
\textbf{2D FFT}
\quad $\rightarrow$ \quad
\textbf{Pointwise multiply}
\quad $\rightarrow$ \quad
\textbf{2D IFFT}

\vspace{10mm}
\large
Instead of repeatedly sliding a kernel, convolution becomes \textbf{element-wise multiplication in the frequency domain}.\textsuperscript{[1]}
\end{frame}

% ------------------------------------------------
\begin{frame}{Frequency-Domain Filtering}
\begin{columns}[T,onlytextwidth]
\column{0.48\textwidth}
\centering
\Large Low-pass
\vspace{2mm}
\begin{itemize}\small
  \item Keeps low spatial frequencies.
  \item Suppresses sharp spatial changes.
  \item Produces a smoother / more blurred heat field.
\end{itemize}

\column{0.48\textwidth}
\centering
\Large High-pass
\vspace{2mm}
\begin{itemize}\small
  \item Removes the DC/background component.
  \item Keeps high spatial frequencies.
  \item Highlights sharp spatial gradients and edges.
\end{itemize}
\end{columns}

\vfill
\begin{center}
\large
\textbf{Heat diffusion itself behaves as a spatial low-pass process.}\textsuperscript{[1]}
\end{center}
\end{frame}

\end{document}
